% =====================================================================
%  figures.tex —— slot-chain-spec 的 10 幅 TikZ 图（由 mermaid 转写）
%  黑白灰学术风格；A4 + 2.5cm 边距，单幅宽度 <= 15cm
%  依赖：tikz（xelatex + ctex）
% =====================================================================
\usetikzlibrary{arrows.meta,positioning,shapes.geometric,fit,calc,chains}

\tikzset{
  % 普通矩形节点
  scnode/.style={draw=black, line width=0.5pt, fill=white, align=flush center,
                 font=\footnotesize, inner sep=3pt, minimum height=7mm},
  % 小号矩形节点
  scsm/.style={scnode, font=\scriptsize, minimum height=6mm},
  % 强调（浅灰）
  scfill/.style={fill=gray!15},
  % 强调（深灰）
  scfillb/.style={fill=gray!30},
  % 判定菱形
  scdec/.style={draw=black, line width=0.5pt, fill=gray!15, diamond, aspect=1.9,
                align=flush center, font=\scriptsize, inner sep=1pt},
  % 注记框（时序图 Note / 旁注）
  scnote/.style={draw=black, line width=0.4pt, fill=gray!15, align=flush center,
                 font=\scriptsize, inner sep=3pt},
  % 虚线边框注记
  scdnote/.style={draw=black, line width=0.4pt, dash pattern=on 1.6pt off 1.4pt,
                  fill=white, align=flush center, font=\scriptsize, inner sep=3pt},
  % 实线箭头
  scarr/.style={draw=black, line width=0.5pt, -{Latex[length=1.8mm,width=1.5mm]}},
  % 虚线箭头（对应 mermaid 的 -.->）
  scdarr/.style={scarr, dash pattern=on 1.6pt off 1.4pt},
  % 双向标注箭头
  scdim/.style={draw=black, line width=0.4pt,
                {Latex[length=1.5mm,width=1.2mm]}-{Latex[length=1.5mm,width=1.2mm]}},
  % 无箭头细线
  scline/.style={draw=black, line width=0.5pt},
  % 生命线
  sclife/.style={draw=black, line width=0.4pt, dash pattern=on 1.4pt off 1.4pt},
  % 边标签
  sclab/.style={font=\scriptsize, align=flush center, inner sep=1.5pt, fill=white},
  sclabl/.style={font=\scriptsize, align=flush left, inner sep=1.5pt, fill=white},
}

% ==== FIGURE 1 ====
\begin{figure}[htbp]
\centering
\begin{tikzpicture}
  \node[scnode, text width=46mm] (B) {构建者（按排班轮换）\\每 slot 出块并签名};
  \node[scnode, text width=46mm, below=11mm of B] (SC)
        {签名链（P2P，未上链）\\签名块本身就是预确认};
  \node[scnode, scfill, text width=46mm, below=13mm of SC] (CAND) {L1 候选批次};
  \node[scnode, text width=46mm, below=9mm of CAND] (WIN)
        {结算窗口 $W_{\mathrm{settle}}$\\窗口内最重的已证明候选胜出};
  \node[scnode, scfill, text width=46mm, below=9mm of WIN] (WF)
        {window-final\\提款/桥接从此状态执行};
  \node[scnode, scfillb, text width=46mm, below=9mm of WF] (LF)
        {L1-final};

  \node[scsm, text width=22mm] (AGG) at ($(CAND.west)+(-3.5,0.9)$) {聚合者\\（不受信任）};
  \node[scsm, text width=22mm] (ANY) at ($(CAND.west)+(-3.5,-1.0)$) {任何人\\（兜底）};

  \draw[scarr] (B) -- node[sclabl, right=1mm] {签名覆盖 \texttt{parent\_hash}} (SC);
  \draw[scarr] (SC) -- node[sclabl, right=1mm, text width=32mm]
        {打包 $+$ 有效性证明，\\一笔 L1 交易} (CAND);
  \draw[scarr] (CAND) -- (WIN);
  \draw[scarr] (WIN) -- node[sclabl, right=1mm] {确定性收盘} (WF);
  \draw[scarr] (WF) -- node[sclabl, right=1mm, text width=32mm]
        {落地交易达 $F_{\mathrm{l1}}$ 深度} (LF);

  \draw[scdarr] (AGG.east) -- node[sclab, pos=0.45, above=0pt] {常态落地} (CAND.west);
  \draw[scdarr] (ANY.east) -- node[sclab, pos=0.45, below=0pt] {越阈后兜底} (CAND.west);
\end{tikzpicture}
\caption{从出块到最终的完整通路：构建者签名产生未上链的签名链，任何人都可以把它连同有效性证明落到 L1 成为候选，结算窗口收盘后依次达到 window-final 与 L1-final；落地环节无特权，聚合者与兜底者可互相替代。}
\label{fig:1}
\end{figure}
% ==== END FIGURE 1 ====

% ==== FIGURE 2 ====
\begin{figure}[htbp]
\centering
\begin{tikzpicture}[font=\scriptsize]
  \node[scsm, text width=17mm] (b1) at (0,0)    {slot $s$\\构建者甲出块};
  \node[scsm, text width=17mm] (b2) at (2.8,0)  {slot $s{+}1$\\构建者乙出块};
  \node[scdnote, scfill, text width=17mm, minimum height=6mm] (gp) at (5.6,0)
        {slot $s{+}2$\\构建者丙掉线\\（缺口）};
  \node[scsm, text width=17mm] (b3) at (8.4,0)  {slot $s{+}3$\\构建者丁出块};
  \node[scsm, text width=17mm] (b4) at (11.2,0) {$\cdots$\\链按秒继续};

  \draw[scarr] (b1) -- (b2);
  \draw[scarr] (b3) -- (b4);
  \draw[scline, dash pattern=on 1.2pt off 1.2pt, gray] (b2.east) -- (gp.west);
  \draw[scline, dash pattern=on 1.2pt off 1.2pt, gray] (gp.east) -- (b3.west);

  \draw[scdarr] (b2.north) to[bend left=26]
        node[sclab, above=0.5mm, text width=62mm]
        {缺口：无需证明、无人受罚；$s{+}3$ 的 \texttt{parent\_hash} 直接指向 $s{+}1$ 的块}
        (b3.north);
\end{tikzpicture}
\caption{单个构建者掉线只在时间线上留下一个 slot 的缺口，下一个构建者以最后可见的块为父继续出块，链不停摆。}
\label{fig:2}
\end{figure}
% ==== END FIGURE 2 ====

% ==== FIGURE 3 ====
\begin{figure}[htbp]
\centering
\begin{tikzpicture}[font=\scriptsize]
  % 时间轴
  \draw[scarr] (0,0) -- (11.6,0);
  \node[font=\scriptsize, anchor=west] at (11.7,0) {};
  \foreach \x in {1.7,6.2,10.7}{\draw[scline] (\x,-0.16) -- (\x,0.16);}

  % 轴下：三个时点
  \node[scnote, text width=34mm, anchor=north] at (1.7,-0.35)
        {快照点 $H_{\mathrm{snap}}(W)$\\$=$ 窗口 $W$ 起点 $-\,D_{\mathrm{snap}}$\\
         取 RANDAO 种子 $+$ 注册表快照\\（已 L1 最终化，不受重组影响）};
  \node[scsm, text width=28mm, anchor=north] at (6.2,-0.35)
        {窗口 $W$ 起点\\slot $= W \times W_{\mathrm{size}}$};
  \node[scsm, text width=20mm, anchor=north] at (10.7,-0.35)
        {窗口 $W$ 终点};

  % 轴上：两段区间
  \draw[scdim] (1.7,0.55) -- (6.2,0.55);
  \node[sclab, text width=44mm, anchor=south] at (3.95,0.65)
        {$D_{\mathrm{snap}} \ge H_{\mathrm{look}} + F_{\mathrm{final}} +$ 余量\\（$=5$ 个 L1 epoch）};
  \draw[scdim] (6.2,0.55) -- (10.7,0.55);
  \node[sclab, text width=42mm, anchor=south] at (8.45,0.65)
        {$W_{\mathrm{size}} = 768$ slot\\固定对齐分区，非滑动窗口};

  % 前瞻区间（虚线）
  \draw[scdarr] (6.2,-2.35) -- (10.7,-2.35);
  \node[sclab, text width=52mm, anchor=north] at (8.45,-2.45)
        {排班在被使用前已定：任一时刻至少 $H_{\mathrm{look}} = 768$ slot 可算};
\end{tikzpicture}
\caption{快照点、排班窗口与前瞻视界的时间几何：熵与注册表在窗口起点前 $D_{\mathrm{snap}}$ 处一次性冻结，因而任一时刻均有至少一个完整窗口的排班可被任何人计算。}
\label{fig:3}
\end{figure}
% ==== END FIGURE 3 ====


% ==== FIGURE 4 ====
\begin{figure}[htbp]
\centering
\begin{tikzpicture}
  \node[scsm, text width=62mm] (S)
        {slot $s$ 的构建者选父块 $P$\\（要求 $P.\mathrm{slot} < s$）};
  \node[scdec, text width=26mm, below=11mm of S] (G)
        {缺口 $s - P.\mathrm{slot}$\\$\le G_{\max}$ ?\\（$G_{\max} = 64$ slot）};
  \node[scsm, scfill, text width=52mm, right=14mm of G] (T1)
        {档 (i) 正常出块\\
         $P$ 只需是签名链中的真实前块\\
         落地状态无关：未落地、已落地、\\
         已 final 均可};
  \node[scsm, scfill, text width=62mm, below=15mm of G] (T2)
        {档 (ii) 恢复（真停摆）\\
         $P$ 必须同时满足：\\
         (1) window-final（所在候选已收盘）\\
         (2) L1-final（落地交易达 $F_{\mathrm{l1}}$ 深度）};
  \node[scsm, text width=62mm, below=7mm of T2] (FR)
        {块头携带 \texttt{final\_ref}（$\le$ 本块 anchor）\\
         证明【签名时】$P$ 已 $F_{\mathrm{l1}}$-final\\
         （堵住 $F_{\mathrm{l1}}$ 边界的预签抢落循环）};
  \node[scsm, text width=62mm, below=7mm of FR] (LAND)
        {落地时另验：$P$ 须是当前 canonical 落地尖端\\
         尖端已被推进 $\Rightarrow$ 恢复块良性搁浅，重建即可};

  \draw[scarr] (S) -- (G);
  \draw[scarr] (G) -- node[sclab, above=0.4mm] {是} (T1);
  \draw[scarr] (G) -- node[sclab, left=1mm, align=flush right] {否\\（档 (i) 不可用）} (T2);
  \draw[scarr] (T2) -- (FR);
  \draw[scarr] (FR) -- (LAND);
\end{tikzpicture}
\caption{两档父块规则的判定流程：判据是签名链上的缺口大小而非父块的落地状态，故签名时、证明时、落地时三处求值一致；只有缺口超过 $G_{\max}$ 的真停摆才进入需要 final 父块的恢复档。}
\label{fig:4}
\end{figure}
% ==== END FIGURE 4 ====

% ==== FIGURE 5 ====
\begin{figure}[htbp]
\centering
\begin{tikzpicture}
  \node[scnode, text width=46mm] (IDLE)
        {无开启窗口\\当前窗口最终头 $= F$};
  \node[scnode, scfill, text width=46mm, below=27mm of IDLE] (OPEN)
        {窗口开启\\best $=$ 当前最重候选\\（provisional，可被取代）};
  \node[scnode, scfillb, text width=46mm, below=14mm of OPEN] (CLOSE)
        {收盘：唯一一次原子提交\\把 best 的终局元组\\写为新 canonical\\
         （推进游标，tip 成为\\新最终头 $F'$）};

  % 主转移（标签置于左侧）
  \draw[scarr] (IDLE) -- node[sclab, left=1.5mm, text width=38mm]
        {首个延伸 $F$ 的合法候选落地：\\
         冻结基线 base $=$\\
         $(F,\ \mathrm{stateRoot},$\\
         $F_{\mathrm{consumed}},\ m_{\mathrm{consumed}})$\\
         \texttt{close\_at} $=$ L1 高度 $+\ W_{\mathrm{settle}}$} (OPEN);
  \draw[scarr] (OPEN) -- node[sclab, left=1.5mm, text width=38mm]
        {L1 高度到达 \texttt{close\_at}\\（同一高度：先收盘、后受理）} (CLOSE);

  % IDLE 虚线自环（右侧）
  \draw[scdarr] ($(IDLE.east)+(0,0.22)$) -- ++(0.8,0) |- ($(IDLE.east)+(0,-0.22)$);
  \node[sclab, text width=36mm, anchor=west] at (3.6,0)
        {收盘后才到的候选\\只能作为下一窗口的首候选};

  % OPEN 自环 1
  \draw[scarr] ($(OPEN.east)+(0,0.46)$) -- ++(0.8,0) |- ($(OPEN.east)+(0,0.16)$);
  \node[sclab, text width=36mm, anchor=west] at ($(OPEN.east)+(1.2,0.78)$)
        {更重合法候选：对同一冻结\\base 验证，key 严格更大\\$\Rightarrow$ 取代 best（只记账）};
  % OPEN 自环 2
  \draw[scarr] ($(OPEN.east)+(0,-0.16)$) -- ++(0.8,0) |- ($(OPEN.east)+(0,-0.46)$);
  \node[sclab, text width=36mm, anchor=west] at ($(OPEN.east)+(1.2,-0.80)$)
        {更差候选到达 $\Rightarrow$ revert，\\零状态变化};

  % 收盘后回到空闲
  \draw[scarr] (CLOSE.west) -- ++(-4.4,0) |- (IDLE.west);
\end{tikzpicture}
\caption{结算窗口的生命周期状态机：窗口开启即冻结基线，窗口内所有候选对同一基线验证且只记账，直到收盘时才把最重候选的终局元组一次性原子提交为新的 canonical 状态。}
\label{fig:5}
\end{figure}
% ==== END FIGURE 5 ====

% ==== FIGURE 6 ====
\begin{figure}[htbp]
\centering
\begin{tikzpicture}
  \node[scnode, text width=48mm] (P1)
        {第 1/2 档：预确认\\构建者签名，秒级};
  \node[scnode, text width=48mm, below=13mm of P1] (P2)
        {provisional 落地\\候选已上 L1，\\窗口内仍可被更重者取代\\
         稳态 $\mathrm{lag}_{\mathrm{prov}} \approx 15$--$20$ min};
  \node[scnode, scfill, text width=48mm, below=13mm of P2] (P3)
        {window-final\\提款/桥接从此执行\\
         稳态 $\mathrm{lag}_{\mathrm{final}} \approx 35$--$40$ min};
  \node[scnode, scfillb, text width=48mm, below=13mm of P3] (P4)
        {L1-final\\绝对最终};

  \draw[scarr] (P1) -- (P2);
  \draw[scarr] (P2) -- node[sclab, left=1.5mm, text width=24mm]
        {$+\ W_{\mathrm{settle}}$\\$\approx 20$ min} (P3);
  \draw[scarr] (P3) -- node[sclab, left=1.5mm, text width=24mm]
        {$+\ F_{\mathrm{l1}}$\\（分钟级）} (P4);

  \node[scdnote, text width=42mm, anchor=west] (N1) at (3.4,-3.0)
        {服务观测阈值\\$\Delta_{\mathrm{lag,prov}} = 4$ epoch $\approx 25.6$ min};
  \node[scdnote, text width=42mm, anchor=west] (N2) at (3.4,-6.1)
        {兜底/计次阈值\\$\Delta_{\mathrm{lag,final}} \approx 9$ epoch $\approx 57.6$ min\\
         $=\Delta_{\mathrm{lag,prov}} + W_{\mathrm{settle,max}} +$ 余量};

  \draw[scdarr] (P2.east) -- (P2.east -| N1.west);
  \draw[scdarr] (P3.east) -- (P3.east -| N2.west);
\end{tikzpicture}
\caption{四级确定性阶梯与两个滞后量：预确认为秒级，provisional 落地与 window-final 之间相差一个结算窗口；$\mathrm{lag}_{\mathrm{prov}}$ 只用于服务节奏观测，$\mathrm{lag}_{\mathrm{final}}$ 才进入兜底与罚则会计。}
\label{fig:6}
\end{figure}
% ==== END FIGURE 6 ====

% ==== FIGURE 7 ====
\begin{figure}[htbp]
\centering
\begin{tikzpicture}
  \node[scnode, text width=52mm] (W)
        {观测 $\mathrm{lag}_{\mathrm{final}}$（L1 状态的纯函数）};
  \node[scdec, text width=26mm, below=10mm of W] (C)
        {$\mathrm{lag}_{\mathrm{final}} > \Delta_{\mathrm{lag,final}}$ ?};
  \node[scsm, text width=30mm, right=12mm of C] (N)
        {兜底窗口关闭\\常态：只有聚合者落地};
  \node[scnode, scfill, text width=52mm, below=12mm of C] (O)
        {兜底窗口开启（状态条件，非计时器）\\
         资格与报酬承诺【快照】保持到结算窗口收盘};
  \node[scnode, text width=52mm, below=8mm of O] (A)
        {任何人可落地合法候选\\
         报酬 $=$ 成本 $+$ 加成，从聚合者保证金支付};
  \node[scdec, text width=28mm, below=12mm of A] (S)
        {开窗期内被接受的\\批次来自谁？};
  \node[scsm, text width=34mm] (K1) at ($(S.south)+(-4.6,-1.5)$)
        {兜底计次 strike\\按 $G_{\mathrm{strike}}$ 限速累积};
  \node[scsm, text width=34mm] (K2) at ($(S.south)+(4.6,-1.5)$)
        {迟到罚金（按超出量定价）\\$+$ 迟到计数（独立限速）};
  \node[scdec, text width=30mm] (T) at ($(S.south)+(0,-3.9)$)
        {累计 $m_{\mathrm{agg}} = 2$ 次\\或 $m'_{\mathrm{agg}} = 4$ 次？};
  \node[scnode, scfillb, text width=44mm, below=8mm of T] (E)
        {席位终止 $\Rightarrow$ 最高候补者晋升};

  \draw[scarr] (W) -- (C);
  \draw[scarr] (C) -- node[sclab, above=0.4mm] {否} (N);
  \draw[scarr] (C) -- node[sclab, left=1mm] {是} (O);
  \draw[scarr] (O) -- (A);
  \draw[scarr] (A) -- (S);
  \draw[scarr] (S.west) -- node[sclab, midway] {非聚合者} (K1.north);
  \draw[scarr] (S.east) -- node[sclab, midway] {聚合者自己（迟到）} (K2.north);
  \draw[scarr] (K1.south) -- (T.west);
  \draw[scarr] (K2.south) -- (T.east);
  \draw[scarr] (T) -- node[sclab, left=1mm] {是} (E);
  \draw[scarr] (T.west) -- ++(-4.8,0)
        node[sclab, above=0.4mm, pos=0.75] {否} |- (W.west);
\end{tikzpicture}
\caption{无许可兜底窗口与席位计次的因果链：兜底资格由 L1 可观测的 $\mathrm{lag}_{\mathrm{final}}$ 越阈这一状态条件开启，开窗期内谁落地决定计入兜底计次还是迟到罚金，计次累积到阈值即触发席位终止。}
\label{fig:7}
\end{figure}
% ==== END FIGURE 7 ====

% ==== FIGURE 8 ====
\begin{figure}[htbp]
\centering
\begin{tikzpicture}[font=\scriptsize]
  \def\pL{0}\def\pO{3.6}\def\pB{7.2}\def\pP{10.6}

  % 参与者
  \node[scnote, text width=22mm, anchor=north] at (\pL,0) {落地者\\（聚合者或任何人）};
  \node[scnote, text width=22mm, anchor=north] at (\pO,0) {L1 合约};
  \node[scnote, text width=22mm, anchor=north] at (\pB,0) {任一回归的\\排班构建者};
  \node[scnote, text width=22mm, anchor=north] at (\pP,0) {P2P 网络};

  % 生命线
  \foreach \x in {\pL,\pO,\pB,\pP}{\draw[sclife] (\x,-1.05) -- (\x,-10.1);}

  % Note over 落地者..P2P
  \node[scnote, text width=104mm] at (5.3,-1.7)
        {停摆发生，缺口 $> G_{\max}$，墙钟远超链头};

  % (1)
  \draw[scarr] (\pL,-3.1) -- (\pO,-3.1);
  \node[sclab, text width=42mm, anchor=south] at (1.8,-3.02)
        {(1) 先把停摆前的最优尾巴\\（含冻结长尾）带证明落为候选};

  % (2) L1 自消息（虚线）
  \draw[scdarr] (\pO,-4.1) -- ++(0.85,0) -- ++(0,-0.55) -- (\pO,-4.65);
  \node[sclab, text width=44mm, anchor=west] at (\pO+1.05,-4.38)
        {(2) 结算窗口收盘 $\Rightarrow$\\尾巴 window-final，最终头 $=F$};

  % (3) 构建者自消息
  \draw[scarr] (\pB,-5.75) -- ++(0.85,0) -- ++(0,-0.55) -- (\pB,-6.30);
  \node[sclab, text width=36mm, anchor=west] at (\pB+1.05,-6.02)
        {(3) 等 $F$ 的落地交易达\\$F_{\mathrm{l1}}$ 深度（分钟级）};

  % (4)
  \draw[scarr] (\pB,-7.6) -- (\pP,-7.6);
  \node[sclab, text width=46mm, anchor=south] at (8.9,-7.52)
        {(4) 以 $F$ 为父签档 (ii) 恢复块\\（\texttt{final\_ref} 见证父已 final）};

  % Note over 构建者..P2P
  \node[scnote, text width=56mm] at (8.9,-8.6)
        {一跳接回墙钟；此后按档 (i) 正常出块，与停摆时长无关};

  % (5)
  \draw[scarr] (\pL,-9.9) -- (\pO,-9.9);
  \node[sclab, text width=54mm, anchor=south] at (1.8,-9.82)
        {(5) 恢复批次照常证明、落地、开新窗口};
\end{tikzpicture}
\caption{数小时乃至数天停摆后的恢复时序：先把停摆前的最优尾巴落地并收盘为最终头 $F$，再由任一回归构建者以 $F$ 为父出一个档 (ii) 恢复块，一跳把排序接回墙钟，恢复过程不含独立子协议。}
\label{fig:8}
\end{figure}
% ==== END FIGURE 8 ====

% ==== FIGURE 9 ====
\begin{figure}[htbp]
\centering
\begin{tikzpicture}
  \node[scnode, scfill, text width=80mm] (Q)
        {块首消费：两条 FIFO 队列，逐块义务};
  \node[scnode, text width=80mm, below=8mm of Q] (BR)
        {桥接队列前缀：\\
         同时满足：条数 $\le C_{\mathrm{bridge}}$\\
         且累计声明 gas $\le C_{\mathrm{bridge}}$ gas 份额\\
         的最长前缀（保留配额）};
  \node[scnode, text width=80mm, below=8mm of BR] (OR)
        {普通强制队列前缀：\\
         条数 $\le C_{\mathrm{force}} - $ 已取桥接条数\\
         累计 gas $\le C_{\mathrm{force}}$ gas 份额 $-$ 桥接已用 gas\\
         的最长前缀};
  \node[scnode, text width=80mm, below=8mm of OR] (R)
        {取到最长前缀即合法，少于它才不合法\\
         （不是「恰好某个条数」）；余量确定性溢出到后续块};
  \node[scnode, scfill, text width=80mm, below=8mm of R] (D)
        {对前置状态非法的条目：弃置消费\\
         游标照常推进、不执行};

  \draw[scarr] (Q) -- (BR);
  \draw[scarr] (BR) -- (OR);
  \draw[scarr] (OR) -- (R);
  \draw[scarr] (R) -- (D);
\end{tikzpicture}
\caption{每个 L2 块块首的强制条目消费规则：桥接队列先按保留配额取最长前缀，普通强制队列再在剩余条数与 gas 配额内取最长前缀，状态非法条目按弃置消费处理以避免队头死锁。}
\label{fig:9}
\end{figure}
% ==== END FIGURE 9 ====

% ==== FIGURE 10 ====
\begin{figure}[htbp]
\centering
\begin{tikzpicture}[font=\scriptsize]
  \draw[scarr] (0,0) -- (10.9,0);
  \foreach \x in {0.9,4.3,7.1,10.3}{\draw[scline] (\x,-0.15) -- (\x,0.15);}

  % 轴下：四个时点
  \node[scsm, text width=30mm, anchor=north] at (0.9,-0.35)
        {块产生\\anchor 已老 $\le D_{\mathrm{anchor}}$\\（32 L1 slot）};
  \node[scsm, text width=20mm, anchor=north] at (4.3,-0.35)
        {兜底被授权};
  \node[scsm, text width=18mm, anchor=north] at (7.1,-0.35)
        {证明完成};
  \node[scsm, scfill, text width=30mm, anchor=north] at (10.3,-0.35)
        {候选落地\\anchor 新鲜度在此校验};

  % 轴上：三段
  \draw[scdim] (0.9,0.5) -- (4.3,0.5);
  \node[sclab, text width=34mm, anchor=south] at (2.6,0.6)
        {聚合者扣落地拖延\\$\le \Delta_{\mathrm{lag,final}}$\\（兜底开窗阈值）};
  \draw[scdim] (4.3,0.5) -- (7.1,0.5);
  \node[sclab, text width=24mm, anchor=south] at (5.7,0.6)
        {证明\\$\le P_{\mathrm{prove,max}}$};
  \draw[scdim] (7.1,0.5) -- (10.3,0.5);
  \node[sclab, text width=28mm, anchor=south] at (8.7,0.6)
        {L1 纳入\\$\le T_{\mathrm{include,max}}$};

  % 全程跨度 + 不变量
  \draw[scdim] (0.9,-2.35) -- (10.3,-2.35);
  \node[scnote, text width=104mm, anchor=north] at (5.6,-2.5)
        {不变量：$D_{\mathrm{anchor,max}} \ge D_{\mathrm{anchor}} + \Delta_{\mathrm{lag,final}}
         + P_{\mathrm{prove,max}} + T_{\mathrm{include,max}} +$ 余量\\
         $\approx 420$ L1 slot $\approx 84$ 分钟，罩住兜底授权全程，死锁不存在};
\end{tikzpicture}
\caption{最坏路径下 anchor 年龄的时间几何：新鲜度窗口 $D_{\mathrm{anchor,max}}$ 必须罩住从块产生、经聚合者拖延至越阈授权兜底、再到证明与 L1 纳入的全程，这一参数关系即为陈旧-anchor 尾巴死锁的解。}
\label{fig:10}
\end{figure}
% ==== END FIGURE 10 ====
